\chapter{Second derivatives}
\label{ch:second-derivatives}

\section{Purpose and scope}

The package differentiates its energy chain twice, so that harmonic
frequencies, curvature checks, and trust-region steps come out of the
same code path that produced the gradients.  The target quantity is the
nuclear Hessian of the relaxed energy $E^{*}(R)$ of
Chapter~\ref{ch:differentiable-chain}, not the Hessian at frozen
orbitals.  The frozen object is available as an intermediate, but it is
not the physical second derivative: the orbital response contributes at
the same order and, on the systems measured here, moves every frequency
in the same direction.

The chapter applies to any energy whose orbitals are determined by a
stationarity condition and whose Hessian is available through
automatic differentiation.  It does not cover analytic second
derivatives of the single-reference SCF surface, and it does not cover
a matrix-free solver that removes the iteration cap of the
least-squares response; both are listed as open problems in the
README.

\section{Theory and derivation}
\label{sec:hessian-theory}

\subsection{Implicit differentiation of the stationarity condition}

Let $g(\theta,R)=\partial E/\partial\theta$ and let
$\theta^{*}(R)$ satisfy
\begin{equation}
  g\bigl(\theta^{*}(R),R\bigr)=0.
  \label{eq:hess-stationarity}
\end{equation}
Differentiating Eq.~\eqref{eq:hess-stationarity} with respect to $R$,
\begin{equation}
  H_{\theta\theta}\frac{d\theta^{*}}{dR}
  + \frac{\partial g}{\partial R}=0,
  \qquad
  H_{\theta\theta}
  = \frac{\partial^{2}E}{\partial\theta^{2}},
  \qquad
  \frac{\partial g}{\partial R}
  = \frac{\partial^{2}E}{\partial\theta\,\partial R}
  \equiv M,
  \label{eq:hess-implicit}
\end{equation}
so
\begin{equation}
  \frac{d\theta^{*}}{dR} = -H_{\theta\theta}^{-1}M.
  \label{eq:hess-dtheta}
\end{equation}
The relaxed energy is $E^{*}(R)=E(\theta^{*}(R),R)$.  Its first
derivative is the envelope result
\begin{equation}
  \frac{dE^{*}}{dR}
  = \left.\frac{\partial E}{\partial R}\right|_{\theta},
  \label{eq:hess-first}
\end{equation}
and its second derivative follows from differentiating
Eq.~\eqref{eq:hess-first}:
\begin{equation}
  \frac{d^{2}E^{*}}{dR^{2}}
  = \left.\frac{\partial^{2}E}{\partial R^{2}}\right|_{\theta}
  + M^{T}\frac{d\theta^{*}}{dR}
  = H_{RR} - M^{T}H_{\theta\theta}^{-1}M.
  \label{eq:hess-total}
\end{equation}
Here $H_{RR}=\partial^{2}E/\partial R^{2}|_{\theta}$ is the frozen
Hessian, and
\begin{equation}
  H_{\mathrm{response}} = -M^{T}H_{\theta\theta}^{-1}M
  \label{eq:hess-response}
\end{equation}
is the orbital-response contribution.  Equation~\eqref{eq:hess-total}
is the construction implemented by
\texttt{core.hessian.implicit\_hessian}; the result object exposes
$H_{RR}$, $H_{\mathrm{response}}$, and the total separately.

\subsection{The pseudo-inverse and null-space exclusion}

The orbital Hessian is not positive definite and is singular along
directions that leave the energy invariant.  Rotations inside a block
are the physical example: the core block is fully occupied, the virtual
block is empty, and the active block is complete inside the CI step, so
a rotation that only mixes orbitals inside one block changes no
occupied--virtual or occupied--active matrix element.  The optimized
$\theta$ is also free to move along any direction that the optimizer
never constrained.

Write the symmetric eigendecomposition
\begin{equation}
  H_{\theta\theta} = \sum_i \lambda_i\,v_i v_i^{T},
  \qquad
  v_i^{T}v_j=\delta_{ij}.
  \label{eq:hess-eig}
\end{equation}
The truncated pseudo-inverse with relative cutoff $\tau$ is
\begin{equation}
  H_{\theta\theta}^{+}
  = \sum_{|\lambda_i|>\tau\lambda_{\max}}
    \lambda_i^{-1}\,v_i v_i^{T},
  \label{eq:hess-pinv}
\end{equation}
and the response becomes
\begin{equation}
  \frac{d\theta^{*}}{dR}
  = -\sum_{|\lambda_i|>\tau\lambda_{\max}}
    \lambda_i^{-1}\,v_i\left(v_i^{T}M\right).
  \label{eq:hess-response-sum}
\end{equation}
A dropped direction contributes to $d\theta^{*}/dR$ only through
$v_i^{T}M$.  If $M$ has no component along the null space, dropping the
direction changes nothing.  If it does, the response along that
direction is not determined by the second-order expansion, and the
pseudo-inverse silently omits it.  The result object exposes
$H_{\theta\theta}$ and $M$, so the projection $v_i^{T}M$ can be formed
for any discarded direction; the curvature report records the spectrum
and the near-null count, not the projection itself.

A relative cutoff is not interchangeable with an absolute one.  An
earlier script used an absolute singular-value threshold of $10^{-3}$;
on a matrix with condition number $10^{12}$ the absolute and relative
rules retained five and two singular values and produced solutions
differing by $99.9\%$.  The default in \texttt{core.hessian} is
relative, \texttt{rcond}$=10^{-9}$.

\subsection{The unrolled route and its consistency}

The alternative construction differentiates the optimizer itself.
Write one Newton-type step as
\begin{equation}
  \theta_{k+1}(R)
  = \theta_k(R) - H_k^{-1}g\bigl(\theta_k(R),R\bigr),
  \label{eq:hess-newton-map}
\end{equation}
and evaluate $N$ steps inside the computational graph.  At a fixed
point $\theta_N=\theta^{*}$ the derivative of the map satisfies
\begin{equation}
  \frac{d\theta_{k+1}}{dR}
  = \left(I-H_k^{-1}H_k\right)\frac{d\theta_k}{dR}
    - H_k^{-1}M_k
  = -H_k^{-1}M_k.
  \label{eq:hess-unrolled}
\end{equation}
With an exact Hessian the fixed-point derivative is reached in one
step; with the quasi-Newton or damped Hessian used in practice the
iteration converges linearly with the contraction factor of the
optimizer.  The unrolled and implicit routes therefore agree in the
limit of a converged iteration and an exact response solve.  The
measured agreement on the validation systems is
$1.6\times10^{-11}$ (Chapter~\ref{ch:verification}).

The unrolled route has two costs that the implicit route avoids.  It
inherits every numerical accident of the optimizer, and its memory and
time scale with the number of unrolled steps.  It is retained as the
independent check, not as the production route.

\subsection{Hessian-vector products}

The refinement and trust-region stages need the Hessian only through
its action on a vector.  The product is the directional derivative of
the gradient,
\begin{equation}
  H v = \left.\frac{d}{d\epsilon}
        g(\theta+\epsilon v)\right|_{\epsilon=0}
      = \mathrm{jvp}_v(g)(\theta),
  \label{eq:hess-hvp}
\end{equation}
and one forward-mode tangent per reverse pass is enough.  Forming the
full Hessian with \texttt{jax.hessian} is forward-over-forward: every
intermediate is broadcast over the square of the parameter count.  At
about five hundred parameters that construction requested a single
$402$~GB allocation and died; the Hessian-vector form is what keeps the
memory at one tangent vector.

\section{Input and operation}
\label{sec:hessian-input}

The production entry point is
\begin{lstlisting}[language=Python]
from diffmrc.core.hessian import (implicit_hessian,
                                  make_joint_energy_fn,
                                  make_theta_energy_fn,
                                  implicit_hessian_columnwise,
                                  implicit_hessian_hvp)
\end{lstlisting}
\texttt{implicit\_hessian(mol, theta, E\_joint, rcond=1e-9)} returns an
\texttt{ImplicitHessianResult} with the fields \texttt{H\_tt},
\texttt{H\_tt\_inverse}, \texttt{curvature}, \texttt{mixed},
\texttt{H\_RR}, \texttt{H\_response}, \texttt{H\_total},
\texttt{dtheta\_dR\_max}, and \texttt{rcond}.  The total is in
Eh/bohr$^{2}$ with shape \texttt{(3 natm, 3 natm)}.  \texttt{theta} is
in radians; \texttt{rcond} is a relative cutoff on the singular values
of \texttt{H\_tt}.

\texttt{E\_joint} is built by \texttt{make\_joint\_energy\_fn} and has
the signature \texttt{(mol, theta) -> energy}.  The SCF is rebuilt
inside every call, which is what makes the geometry derivative correct
and what makes each evaluation expensive.  \texttt{E\_tt}, when
available from \texttt{make\_theta\_energy\_fn}, removes the SCF from
the theta-only part of the graph; the mixed and frozen blocks still use
the joint function because they differentiate the reference orbitals
with respect to geometry.

For large parameter counts, \texttt{implicit\_hessian\_columnwise}
takes one tangent per column, and \texttt{implicit\_hessian\_hvp}
assembles the total without forming \texttt{H\_tt}, \texttt{M}, or
\texttt{H\_RR}.  The matrix-free route solves
\texttt{H\_tt x = M v} by LSMR and records \texttt{lsmr\_iterations},
\texttt{lsmr\_istop}, and \texttt{lsmr\_residual} per column.  A
converged solve is part of the result; the iteration cap is
$8n_{\mathrm{par}}$.

Harmonic frequencies are computed from the total Hessian with the
\texttt{wavenumbers} routine in \texttt{analysis.vibrations}, which
performs the mass-weighting and the normal-mode analysis.  The Hessian must be
handed over in the (atom, coordinate) block layout the thermodynamics
module expects.  A plain reshape compiles and returns a plausible
number; for CO/6-31G it produced $1501.8$~cm$^{-1}$ against the correct
$2293.5$~cm$^{-1}$.

\section{Worked example}
\label{sec:hessian-example}

\subsection{H\texorpdfstring{$_2$}{2}/STO-3G: orbital-space HVP against
a finite-difference ladder}

The HVP test in \texttt{tests/driver/test\_nevpt2.py} uses the
NEVPT2 joint energy on H$_2$/STO-3G, a symmetric bond-stretch probe
vector, and a ladder of central differences of the AD gradient.  The
recorded errors of the HVP against the ladder are
\begin{center}
\small
\begin{tabular}{lr}
\toprule
Step $\delta$ (bohr) & $\max|\mathrm{HVP}-\mathrm{FD}|$ \\
\midrule
$1\times10^{-3}$ & $6.6\times10^{-6}$ \\
$3\times10^{-4}$ & $6.0\times10^{-7}$ \\
\bottomrule
\end{tabular}
\end{center}
The ratio is approximately $(10/3)^2$, which is the expected quadratic
truncation term of the central difference.  The test requires monotone
improvement and the bounds $2\times10^{-5}$ at $\delta=10^{-3}$ and
$2\times10^{-6}$ at $\delta=3\times10^{-4}$.  This is an HVP in the
orbital-parameter space, not in the nuclear-coordinate space; the
probe vector is the bond stretch because a generic displacement of a
diatomic would compare zero against zero.

\subsection{H\texorpdfstring{$_2$}{2}O/6-31G: relaxed Hessian against
the golden finite-difference standard}

The recorded H$_2$O/6-31G CAS(2,2) Hessian diagonal for the three
oxygen displacements is
\begin{center}
\small
\begin{tabular}{lrrr}
\toprule
Route & $H_{xx}$ & $H_{yy}$ & $H_{zz}$ \\
\midrule
golden standard (finite differences) & $0.7192711493$ & $0.0217924736$ & $0.4702696863$ \\
unrolled iteration, $N=3$ & $0.7192706205$ & $0.0217921326$ & $0.4702703066$ \\
frozen orbitals (no response) & $0.7261828915$ & $0.0217921326$ & $0.4711140468$ \\
\bottomrule
\end{tabular}
\end{center}
{\raggedright
The values are in Eh/bohr$^{2}$ and come from the \texttt{REF['H2O']}
record in \texttt{experiments/orbital\_response\_implicit.py}; the
golden standard is the Newton-refined reference recorded in the same
file and in \texttt{experiments/orbital\_response\_generalize.py}.  The
unrolled route differs from the golden standard by
$5.3\times10^{-7}$, $3.4\times10^{-7}$, and $6.2\times10^{-7}$~Eh/bohr$^2$
(relative $7.4\times10^{-7}$, $1.6\times10^{-5}$, and
$1.3\times10^{-6}$).  The frozen route misses the first and third
components by $6.9\times10^{-3}$ and $8.4\times10^{-4}$~Eh/bohr$^2$:
the response term is not a small correction.  The implicit route
reproduces the unrolled construction to $1.6\times10^{-11}$.\par}

\subsection{CO/6-31G: the frequency anchor}

The frequency conversion is anchored against an independent value: CO
at 6-31G gives $2294.2$ and $2293.5$~cm$^{-1}$ from the two assembly
routes, against an external anchor at $2293.5$~cm$^{-1}$.  The
near-agreement of the two routes is the internal check; the external
anchor is the physical check.  A validation target near zero should not
be checked by finite differences: this project measured
finite-difference errors of over one hundredfold on a near-zero mixed
derivative, growing as the step shrank.

\section{Pitfalls}
\label{sec:hessian-pitfalls}

\begin{itemize}
\item \emph{A plain solve on a singular Hessian.}  The orbital Hessian
  reaches a condition number of $2.6\times10^{15}$ on N$_2$.  A direct
  solve inverts the near-null direction and amplifies whatever the
  right-hand side has there.  Use the pseudo-inverse, and report the
  cutoff.
\item \emph{An absolute instead of a relative cutoff.}  The two rules
  retain different subspaces on an ill-conditioned matrix; the measured
  difference was $99.9\%$ in the response vector.  \texttt{rcond} is
  relative.
\item \emph{Forgetting \texttt{stop\_gradient} on $\theta$ in the
  frozen Hessian.}  The frozen term then silently acquires the orbital
  response, and the two routes stop being distinguishable.
\item \emph{Unsymmetrized AD Hessians.}  AD returns a matrix symmetric
  only to round-off; \texttt{eigvalsh} and \texttt{pinv} assume exact
  symmetry.  Symmetrize as $0.5(H+H^{T})$.
\item \emph{Finite differences near a zero derivative.}  The measured
  failure is over one hundredfold, with the error growing as the step
  shrinks.  Use an invariance or an analytic cross-check.
\item \emph{The frequency layout trap.}  A plain reshape of the Hessian
  gives a plausible but wrong spectrum ($1501.8$ against
  $2293.5$~cm$^{-1}$ for CO/6-31G).  Use the (atom, coordinate) block
  layout.
\item \emph{Unblocked eigensolvers and degenerate pairs.}  A plain
  eigensolver's second derivative is NaN at a cross-$S_z$ degeneracy;
  linear molecules need the blocked solver and, for bit-exact
  degenerate pairs, the Lanczos kernel.
\item \emph{Unrolled iteration cost.}  The unrolled route scales with
  the number of steps and inherits optimizer noise; it is the
  independent check, not the production path.
\end{itemize}

\section{Validation and references}
\label{sec:hessian-validation}

{\raggedright
The records behind this chapter are
\texttt{tests/core/test\_hessian.py} (implicit assembly on analytic
surrogates, pseudo-inverse behaviour),
\texttt{tests/core/test\_stationary.py} (Newton--CG and saddle checks),
\texttt{tests/driver/test\_orbitals.py} (Hessian-vector products
against the explicit Hessian), and
\texttt{tests/driver/test\_nevpt2.py} (the H$_2$/STO-3G HVP ladder).
The experimental records are
\texttt{experiments/orbital\_response\_implicit.py},
\texttt{experiments/orbital\_response\_unrolled.py},
\texttt{experiments/orbital\_response\_generalize.py},
\texttt{experiments/hessian\_to\_frequencies.py},
\texttt{experiments/verify\_nevpt2\_an.py}, and
\texttt{experiments/verify\_rcond\_choice.py}.  Chapter~\ref{ch:verification}
collects the measured agreement values: $1.6\times10^{-11}$ between
the implicit and unrolled routes, $10^{-10}$ for the HVP against the
explicit Hessian on a non-quadratic fixture, and the CO frequency
anchor.\par}

The derivation is the implicit-function construction of analytic
second derivatives for MCSCF.  The relevant references are
P.~Pulay, \emph{Mol. Phys.} \textbf{17}, 197 (1969); N.~C.~Handy and
H.~F.~Schaefer, \emph{J. Chem. Phys.} \textbf{81}, 5031 (1984);
T.~Helgaker, P.~J{\o}rgensen and J.~Olsen, \emph{Molecular
Electronic-Structure Theory} (Wiley, 2000), Chapters 10--12; and
J.~A.~Pople, R.~Krishnan, H.~B.~Schlegel and J.~S.~Binkley,
\emph{Int. J. Quantum Chem. Symp.} \textbf{13}, 225 (1979), for the
coupled-perturbed second-derivative construction.
